We next consider the existence of a primitive. Of course, we cannot define a definite integral of a generalized function. Nevertheless, primitives can be shown to exist.

\begin{Theorem}
\label{thm:distribution-primitive-existence}
Let $I=(a,b)$ be an open interval in $\mathbb{R}$ and $f\in\mathcal{D}^{\prime}(I)$. Then there exists $u\in\mathcal{D}^{\prime}(I)$ such that $u'=f$. The primitive $u$ is unique up to a constant.
\end{Theorem}

\begin{Proof}
The uniqueness part is clear from the previous theorem. To construct a primitive, we use the decomposition \eqref{eq:test-function-one-dimensional-zero-integral-decomposition}
\begin{equation}\label{eq:primitive-construction-test-decomposition}
\phi(x)=\phi_1(x)\int_I\phi(s)\,\mathrm{d}s+\psi(x),
\end{equation}
and we let
\begin{equation}\label{eq:primitive-construction-zero-integral-test-primitive}
\chi(x)=\int_a^x\psi(y)\,\mathrm{d}y.
\end{equation}
We then define
\begin{equation}\label{eq:primitive-distribution-definition-pairing}
(u,\phi)=C\int_I\phi(x)\,\mathrm{d}x-(f,\chi),
\end{equation}
where $C$ is an arbitrary constant. If $\phi=\eta'$, then $\int_I\phi(x)\,\mathrm{d}x=0$ and $\psi=\phi$; hence $\chi=\eta$. We thus find
\begin{equation}\label{eq:primitive-distribution-derivative-check}
(u,\eta')=-(f,\eta);
\end{equation}
hence $u'=f$.
\end{Proof}

The multidimensional result that any curl-free vectorfield on a simply connected domain is a gradient can also be extended to distributions; the proof is considerably more complicated than in the one-dimensional case and will not be given here.

The most elementary technique of solving an ODE is based on reducing it to the form $y'=f$; this is why solving an ODE is referred to as "integrating" it. Such procedures also work for distributional solutions. Consider, for example, the ODE
\begin{equation}\label{eq:distribution-linear-first-order-ode}
y'=a(x)y+f(x).
\end{equation}


We assume that $a\in C^\infty(\mathbb{R})$ and $f\in\mathcal{D}^{\prime}(\mathbb{R})$. We can now set
\begin{equation}\label{eq:distribution-ode-integrating-factor-substitution}
y(x)=z(x)\exp\left(\int_0^xa(s)\,\mathrm{d}s\right);
\end{equation}
note that multiplication of distributions by $C^\infty$ functions is well defined and the product rule of differentiation is easily shown to hold. We thus obtain the new ODE
\begin{equation}\label{eq:distribution-ode-primitive-reduction}
z'(x)=f(x)\exp\left(-\int_0^xa(s)\,\mathrm{d}s\right).
\end{equation}


From Theorem~\ref{thm:distribution-primitive-existence}, we know that this ODE has a one-parameter family of solutions.

In particular, if $f$ is a continuous function, then all distributional solutions of \eqref{eq:distribution-linear-first-order-ode} are the classical ones. This is not necessarily true for singular ODEs; for example both the constant 1 and the Heaviside function are solutions of $xy'=0$.
